\section{Discussion / Limitations}
\label{sec:discussion}

A capability ceiling on a cluster of hard motions (crawl/kneel/squat, single-leg dynamic balance)
persists across shape, reward-relaxation, architecture, and morphology-representation
interventions. What did not work and why: FiLM-style conditioning (failed outright, traced to a
12{,}288-output conditioner bottleneck and multiplicative-gradient instability); residual adapters
(v1--v6, plateaued 78--82\%); MoE routing (no separation from a capacity-matched control);
self-attention temporal encoding (no separation from flat-concat; the shared \texttt{Transformer}
module also lacks positional encoding — flagged, deliberately not fixed, judged unlikely to move
the ceiling); physics-derived morphology features (no statistically significant gain over raw
betas, though qualitatively suggestive).

The broader lesson from the design-evolution narrative
(Sections~\ref{sec:exp:architecture}--\ref{sec:exp:morphology}) is that the hard-motion ceiling
looks like it is about \emph{motion difficulty}, not about which architecture or which
conditioning representation carries the signal — every lever tried moved metrics by less than
noise. Resource constraints limited ablation breadth to reduced-scale corpora, addressed directly
in the framing of those sections rather than hidden.

The shape-invariance evaluation (Section~\ref{sec:evaluation}) is a separate axis from this
ceiling discussion: the ceiling is about which \emph{motions} are hard, not about \emph{which
bodies} execute them — consistent with the near-zero shape-failure correlation
(Section~\ref{sec:exp:failure}).
